% Part: lambda-calculus
% Chapter: introduction
% Section: composition

\documentclass[../../../include/open-logic-section]{subfiles}

\begin{document}

\olfileid{lam}{int}{com}
\olsection{\usetoken{S}{lambda definable} تابعې د ترکيب تر عمل لاندې تړلې دي}

\begin{lem}
!!{lambda definable} تابعې د ترکيب تر عمل لاندې تړلې دي۔
\end{lem}

\begin{proof}
فرض کړئ $f$ د $h$، $g_0,$
\dots،~$g_{k-1}$ له ترکيب څخه تعريف شوې ده۔ که $h$، $g_0$،
\dots،~$g_{k-1}$ په ترتيب سره د $H$، $G_0$،
\dots،~$G_{k-1}$ په وسيله !!{lambda defined} وي، بايد داسې
ترم~$F$ ومومو چې~$f$ !!{lambda define}s۔ خو~$F$ په ساده ډول داسې
تعريفولے شو:
\[
F(x_0, \dots, x_{l-1}) = H(G_0(x_0, \dots, x_{l-1}),
\dots, G_{k-1}(x_0, \dots, x_{l-1})).
\]
په بله وينا، د لامبډا حساب ژبه د ترکيب د تمثيل دپاره مناسبه ده۔
\end{proof}

\end{document}
